\section*{Optimal root-KL payment on the positive-residual critical class}

All logarithms are base two.  Let $f:\{-1,1\}^n\to\{-1,1\}$ be
Boolean and write
\[
 p(S)=\widehat f(S)^2,\qquad
 U(f)=\sum_Sp(S)\log_2\frac{\pi(S)}{p(S)},\qquad
 I(f)=\sum_S|S|p(S).
\]
For a physical coordinate $i$, use the signed restriction notation
\[
 a_A=\widehat f(A),\quad b_A=\widehat f(A\cup\{i\}),\quad
 m_A=a_A^2+b_A^2,
\]
\[
 C_A=\sum_{T\supseteq A}a_Tb_T,\quad
 M_A=\sum_{T\supseteq A}m_T,\quad
 \alpha_A=\frac{2a_Ab_A}{m_A},\quad
 \beta_A=\frac{2C_A}{M_A}.
\]
The exact Boolean restriction identity is
\[
 \Delta_iU=B_i-\Gamma_i+K_i,
 \qquad
 K_i=\sum_A m_AD_2\left(\frac{1+\alpha_A}{2}\middle\|
                         \frac{1+\beta_A}{2}\right).       \tag{1}
\]
Boolean convolution gives $C_\varnothing=0$, so the root row of (1) is
\[
 K_i^{(0)}=m\left[1-h_2\left(\frac{1+2ab/m}{2}\right)\right],
 \quad a=\widehat f(\varnothing),\ b=\widehat f(\{i\}),\
 m=a^2+b^2.                                             \tag{2}
\]

\subsection*{1. The sharp scalar coefficient}

\paragraph{Lemma.}
For every real $a,b$,
\[
 \boxed{K_i^{(0)}\le\frac{2}{\ln2}\min\{a^2,b^2\}.}       \tag{3}
\]
The constant $2/\ln2$ is optimal.

\paragraph{Proof.}
By symmetry assume $b^2\le a^2$ and put
\[
 r=\frac{b^2}{a^2+b^2}\le\frac12,
 \qquad x=\frac{2|ab|}{a^2+b^2}=2\sqrt{r(1-r)}.
\]
In natural units define
\[
 d(x)=\frac{1+x}{2}\ln(1+x)+\frac{1-x}{2}\ln(1-x).
\]
Then $d'(x)=\operatorname{artanh}x$.  Moreover
\[
 q(x):=\frac{x}{\sqrt{1-x^2}}-\operatorname{artanh}x
\]
satisfies $q(0)=0$ and
\[
 q'(x)=\frac{1-\sqrt{1-x^2}}{(1-x^2)^{3/2}}\ge0.
\]
Therefore $d'(x)\le x/\sqrt{1-x^2}$, and integration from zero gives
\[
 d(x)\le1-\sqrt{1-x^2}=2r.                            \tag{4}
\]
Since
\[
 1-h_2\left(\frac{1+x}{2}\right)=\frac{d(x)}{\ln2},
\]
equations (2)--(4) prove
\[
 K_i^{(0)}\le\frac{2m r}{\ln2}=\frac{2b^2}{\ln2}.
\]
Interchanging $a,b$ proves (3).  As $|b/a|\to0$,
$d(x)=x^2/2+O(x^4)$ and $x=2|b/a|+O(|b/a|^3)$, so
$K_i^{(0)}/b^2\to2/\ln2$.  Hence the coefficient is sharp. \qed

Define the bias-sensitive singleton reserve
\[
 M_1(f)=\sum_i\min\{p(\varnothing),p(\{i\})\},             \tag{5}
\]
and the nonroot signed debt
\[
 D_+(f)=\sum_i(B_i-\Gamma_i)+
 \sum_i\sum_{A\ne\varnothing}m_{i,A}D_2\left(
 \frac{1+\alpha_{i,A}}2\middle\|\frac{1+\beta_{i,A}}2
 \right).                                                \tag{6}
\]
Summing (1) and applying (3) yields the dimension-free exact consequence
\[
 \boxed{G(f)\le D_+(f)+\frac{2}{\ln2}M_1(f).}              \tag{7}
\]

\subsection*{2. Positive-residual extremal reduction}

Freeze $(a,b)=(4,3/4)$ and set
\[
 R(f)=G(f)-4I(f)-\frac34U(f).                              \tag{8}
\]
The complete residual is additive under disjoint Boolean tensor products.
Thus if the frozen GU inequality fails, a positive-residual tensor-prime
factor exists; equivalently one may choose a dimension-minimal
positive-residual counterexample, which is tensor-prime.  No optimizer over
all dimensions is assumed.

It is therefore sufficient to prove only on tensor-prime Boolean functions
with $R(f)>0$ the aggregate inequality
\[
 \boxed{
 D_+(f)+\frac{2}{\ln2}M_1(f)
 \le4I(f)+\frac34U(f).
 }                                                        \tag{E}
\]
Indeed (7) and (E) imply $R(f)\le0$, contradicting the admitted class.
Tensor additivity then gives
\[
 G(f)\le4I(f)+\frac34U(f)                                 \tag{9}
\]
for every Boolean $f$.  The source reduction consequently gives
\[
 U(f)\le16I(f),\qquad H(f)\le34I(f).                      \tag{10}
\]

This is strictly weaker than the preceding coefficient-four root payment.
That payment consumed the whole $4I$ term and left the requirement
$D_+\le3U/4$.  The new endpoint permits
\[
 D_+\le4I+\frac34U-\frac{2}{\ln2}M_1.                    \tag{11}
\]
Since $M_1\le I$, its right side exceeds $3U/4$ by at least
$(4-2/\ln2)I=1.114609918\ldots I$ for every nonconstant source.
It also retains the exact bias information in (5) rather than replacing it
by singleton influence.

\subsection*{3. Scope and controls}

The even quadratic phases that refute every unqualified fixed
good-coordinate curvature bound have $R<0$ and are excluded at the exact
extremality line (8), not by an added exception.  The new proof changes
mechanism from physical curvature to the signed Boolean root normalization.

The companion computation reconstructs (1), (5), and (6) on all $32,906$
Boolean signs modulo global output sign through dimension four.  It checks
(E) there and the scalar inequality independently.  This is finite evidence
only.  The all-dimensional nonempty-row inequality (E) remains open, as do
universal GU and full FEI.  No claim is made that positive residual is
attained or that descendants inherit criticality.
